\subsection{系统边界与粒子携带的能量}
\label{th:openboundary}
对开口容器应用第一定律时，取固定区域内的瞬时存量作为研究对象。设容器中有 $N$ 个粒子，其内部状态可由 $N,T$ 描述，内能为 $U(N,T)$，短时间内逸出粒子数为 $\dd N_{\rm out}>0$，每粒子平均携带能量为 $\varepsilon_{\rm out}$，则容器粒子数变化为 $\dd N=-\dd N_{\rm out}$。若外界还供热 $\delta Q$，边界无宏观机械位移，
\[
 \dd U=\delta Q-\varepsilon_{\rm out}\dd N_{\rm out}
 =\delta Q+\varepsilon_{\rm out}\dd N.
\]
因 $\dd N<0$，后一项为负。对理想独立粒子，$U=Nu(T)$，完整微分为
\[
 \dd U=Nc_v(T)\dd T+u(T)\dd N.
\]
其中 $u\,\dd N$ 表示存量减少本身造成的内能变化，$Nc_v\dd T$ 表示余气温度改变。与边界输运式联立，
\[
 Nc_v(T)\dd T
 =\delta Q+\bigl[\varepsilon_{\rm out}-u(T)\bigr]\dd N.
\]
余气升温或冷却由逸出能量相对于体内平均能量的高低决定。能量零点若平移 $E_\ast$，$u$ 和 $\varepsilon_{\rm out}$ 同时增加 $E_\ast$，差值不变，温度方程也不变。

\Diagram{fig:controlboundary}{固定边界内的存量与穿越边界的能量}{
 \draw[curve](0,0)--(0,2)--(3.4,2)--(3.4,1.15);
 \draw[curve](3.4,.85)--(3.4,0)--(0,0);
 \node at(1.7,1){$U(N,T)$};
 \draw[axis](-1.6,1)--(0,1)node[midway,above]{$\delta Q$};
 \draw[axis](3.4,1)--(5.6,1)node[right]{$\varepsilon_{\rm out}\dd N_{\rm out}$};
 \node at(2,-.55){先选边界，再写储能与输运};
}
图\ref{fig:controlboundary}中，边界能流按穿越的粒子统计。小孔分子流按穿孔事件取样，平动部分为 $2\kb T$，内部自由度若与穿孔概率无关，则仍取体内平均值。这一微观输运模型与有整体流速的连续稳流模型所用能量形式不同。

对一段连续流管，把入口前的流体推入控制体，外界除送入该流体本身的内能外，还需克服入口压强做功。质量 $\dd m$ 占据体积 $\dd m/\rho$，推入功为 $p\,\dd m/\rho$。因此入口携带的内能加流动功合成比焓 $h=u_m+p/\rho$。在稳定流动中，
\[
 \dot Q-\dot W_{\rm shaft}
 =\dot m\left[
 h_2-h_1+\frac{V_2^2-V_1^2}{2}
 +g(z_2-z_1)\right].
\]
这里 $u_m,h$ 均为单位质量的量，$V_i$ 为宏观流速；$\dot W_{\rm shaft}$ 是对外输出的轴功。入口、出口流动功已经计入焓，不能再作为额外项重复加入。分子流小孔则从单粒子动能与内部能量直接计算输运率；若已这样计入穿面能量，再机械地附加一个 $pV$ 项，便改变了模型。

\subsection{流出冷却与稳流节流的不同约束}
刚性容器缓慢泄流时，容器存量随时间改变，能量方程中有 $\dd U/\dd t$。绝热节流则通常选取阀门附近的固定控制体，在稳态下控制体储能不变；无轴功且宏观动能、位能变化可忽略时，得到 $h_1=h_2$。自由膨胀又是封闭气体占据更大可达体积，外边界刚性绝热，给出整个气体的 $U_1=U_2$。三者分别固定存量能量的变化式、入口出口焓差和整体内能差。

这些约束须在状态方程之前确定。理想气体自由膨胀或理想稳流节流的初末温度相同，是因为它的内能或焓只依赖温度；绝热泄流容器却因穿孔样本偏向高能粒子而冷却。对实际气体，等内能与等焓曲线一般不同，两种温变也须分别计算。

\subsection{持续充气中的流动功}
\begin{example}[绝热刚性空罐的充气温度]
\label{ex:tankfilling}
一个初始接近真空、体积为 $V$ 的刚性容器，通过阀门与温度 $T_s$、压强 $p_s$ 恒定的大型气源相连。气体为常热容理想气体，摩尔热容为 $c_{V,m}$、$c_{p,m}$，$\gamma=c_{p,m}/c_{V,m}$。容器与外界绝热，忽略罐壁热容和最终宏观动能。充气后关阀，待罐内达到均匀平衡。求终温，并说明气源做的功进入了哪里。
\Diagram{fig:tankfillingq}{大型气源经阀门向绝热刚性罐充气}{
 \node[draw,minimum width=2.4cm,minimum height=1.6cm](s)at(0,1){$p_s,T_s$};
 \draw[curve](3,0)rectangle(6.4,2);
 \node at(4.7,1){$V,\ \nu(t),T(t)$};
 \draw[axis](1.2,1)--(3,1);
 \draw(2,.75)--(2.4,1.25);\draw(2,1.25)--(2.4,.75);
 \node[below]at(2.2,.5){阀};
 \node at(4.7,-.5){外壁不动，$Q=0$};
}
\end{example}
\begin{solution}
在图\ref{fig:tankfillingq}中，把连接段和容器共同纳入控制体，在气源侧选入口。每进入 $\dd\nu$ 摩尔气体，气源送入摩尔内能 $c_{V,m}T_s$，并通过入口压强做流动功 $RT_s$，合计送入摩尔焓 $c_{p,m}T_s$。若所选入口处流速不可忽略，应计其动能；本题可在大型气源内部流速可忽略处取截面，直接使用源的摩尔焓。
\[
 \dd U=c_{p,m}T_s\,\dd\nu.
\]
从初始空罐积分，$U_f-U_i=c_{p,m}T_s\int_0^{\nu_f}\dd\nu=c_{p,m}T_s[\nu]_0^{\nu_f}=\nu_fc_{p,m}T_s$，其中$U_i=0$。气体最终静止，内能为 $\nu_fc_{V,m}T_f$，于是
\[
 T_f=\gamma T_s.
\]
\Diagram{fig:tankfillinga}{进入罐内的焓分解为原有内能与入口流动功}{
 \node[draw](a)at(0,0){$c_{p,m}T_s$};
 \node[draw](b)at(4,1){$c_{V,m}T_s$};
 \node[draw](c)at(4,-1){$RT_s$};
 \draw[axis](a)--(b);\draw[axis](a)--(c);
 \node[right]at(6,0){均进入最终内能};
}
图\ref{fig:tankfillinga}中，流动功经过射流、黏滞耗散和内部混合，最终转成罐内热运动能。外壁不动只排除了外壁体积功，入口处的压强仍在输送能量。若充气持续到罐压等于 $p_s$，则 $\nu_f=p_sV/(R\gamma T_s)$。

若罐壁初温为 $T_{w0}$、热容为 $C_w$，且最终与气体同温，则同一控制体给出
\[
 (\nu_fc_{V,m}+C_w)T_f
 =\nu_fc_{p,m}T_s+C_wT_{w0}.
\]
这一步只改变储能项；气源输入的摩尔焓不变。真实充气过程不可逆，均匀终态可用状态方程描述，充气途中则无须假定罐内始终处于准静态。
\end{solution}

